\chapter{Requirements}
\label{chapter:requirements}

This chapter defines the requirements necessary to establish a productivity benchmarking system compliant with the \ac{GDPR} and \ac{BDSG}. Following standard IEEE definitions, requirements are divided into functional, defining what the product must do, and nonfunctional, specifying how the system should operate \autocite{IEEE}. Furthermore, these requirements are structured across three main categories: legality, system architecture, and integration. Specifically, the service must abide by relevant data protection laws, adhere to author-specified system design constraints, and satisfy the integration demands outlined in the MECOIS contributing guidelines \autocite{github_mecois}.

\section{Functional Requirements}
\label{sec:functional_requirements}

The Table \ref{tab:functional_requirements} specifies what the anonymization service must do.
The requirements are sorted by \ac{FLR}, \ac{FSR}, and \ac{FIR}.

\begin{table}[ht]
	\caption{Functional requirements}
	\label{tab:functional_requirements}
	\begin{tabular}{|p{1.5cm}|p{5cm}|p{7cm}|}
		\hline
		ID & Requirement & Description \\
		\hline
		FLR-1 & Data Anonymization & 
		The service must follow the German privacy laws (\ac{GDPR}) and \ac{BDSG}). \\
		\hline
		FSR-1 & Preserve Data & 
		The service must preserve all available Data. \\
		\hline
		FSR-2 & Independent Deployment & 
		The service can be deployed and run independent from the rest of the MECOIS project. \\
		\hline
		FSR-3 & Deterministic Pseudonym Generation & 
		The service must support consistent \ac{UUID} generation across pipeline runs so that cross-platform metrics can be linked over time without revealing identity. \\
		\hline
		FIR-1 & Seamless Integration &
		The service must integrate into the project with as few changes as possible to existing features. \\
		\hline
	\end{tabular}
\end{table}

The requirement for Data Anonymization (FLR-1) is intentionally formulated with a broad scope, as the \ac{GDPR} does not prescribe a singular technical methodology but rather establishes a high-level legal standard. Because various technical approaches can achieve compliance, Chapter \ref{chapter:design_and_implementation} systematically evaluates and compares different anonymization strategies suitable for high-dimensional development metadata. This legal mandate exists in fundamental tension with the system requirement to Preserve Data (FSR-1), which seeks to maintain the integrity of all source information for analytical depth. Consequently, a primary objective of this work is to navigate this conflict by identifying a functional compromise that ensures robust individual privacy while maintaining the data utility required for the derivation of the MECOIS Productivity Index.

\section{Nonfunctional Requirements}
\label{sec:nonfunctional_requirements}

The Table \ref{tab:nonfunctional_requirements} specifies how the anonymization service must operate.
They are sorted by \ac{NFLR}, \ac{NFSR}, and \ac{NFIR}.

 \begin{table}[ht]
 	\caption{Nonfunctional requirements}
 	\label{tab:nonfunctional_requirements}
 	\begin{tabular}{|p{1.5cm}|p{4cm}|p{8cm}|}
 		\hline
 		ID & Requirement & Description \\
 		\hline
 		NFLR-1 & Anonymization Methodology Adherence & 
 		The service must implement the anonymization measures decided upon in Chapter \ref{chapter:design_and_implementation}.  \\
 		\hline
 		NFSR-1 & Containerized Service & 
 		The service must run in a docker container, a common used framework for containerized deployment. \autocite[9--11]{jangla2018accelerating} \\
 		\hline
 		NFSR-2 & RESTful \ac{API} &
 		The service must implement a RESTful \ac{API} to ensure a standardized, stateless, and scalable interface for communication between the Main System and Identity Management System. \\
 		\hline
 		NFIR-1 & Data Integrity & 
 		The service must maintain schema consistency during the serialization and deserialization of DataFrames into \ac{JSON} to ensure no data loss occurs during the \ac{HTTP} transmission process. \\
 		\hline		
 		NFIR-2 & Programming Language & 
 		The Backend of the MECOIS project is written in python3. Therefore, this project must  also use python3. \autocite{github_mecois} \\
 		\hline
 		NFIR-3 & License Policy & 
 		The contributing guidelines state that dependencies should be kept to a minimum and forbid packages with a copyleft or proprietary license.  \autocite{github_mecois} \\
 		\hline
 		NFIR-4 & Styleguide and Clean Code & 
 		Implementations must follow the style guidelines of the main project. \autocite{github_mecois} \\
 		\hline
 	\end{tabular}
 \end{table}